\section{Колебания}

\subsection{Малые гармонические колебания}
Колебания называются \textbf{гармоническими}, если описывающая их величина изменяется со временем по закону синуса или косинуса:
\[
x(t) = A \cos(\omega_0 t + \alpha)
\]

где $A$ — амплитуда, $\omega_0$ — циклическая частота, $\alpha$ — начальная фаза. 

Для \textbf{малых колебаний} при \( \theta\to 0 \) можно линеаризовать уравнения движения из разложения функций в ряд Тейлора в окрестности нуля:

\begin{itemize}
    \item Для синуса:
    \[
    \sin \theta = \theta - \frac{\theta^3}{3!} + \frac{\theta^5}{5!} - \frac{\theta^7}{7!} + \dots
    \]
    При малых \( \theta \) слагаемыми высшего порядка можно пренебречь:
    \[
    \sin \theta \approx \theta
    \]

    \item Для косинуса:
    \[
    \cos \theta = 1 - \frac{\theta^2}{2!} + \frac{\theta^4}{4!} - \frac{\theta^6}{6!} + \dots
    \]
    При малых $\theta$ приближение до первого нетривиального порядка:
    \[
    \cos \theta \approx 1 - \frac{\theta^2}{2}
    \]
\end{itemize}

\subsection{Энергия математического маятника}

\textbf{Математический маятник} представляет собой материальную точку массы \( m \), подвешенную на невесомой и нерастяжимой нити длины \( l \). 

\begin{itemize}
    \item \textbf{Кинетическая энергия:} \( T = \frac{mv^2}{2} \), где скорость \( v = l\dot{\theta} \), \( \theta \) — угол отклонения от вертикали.
    \[
    T = \frac{m}{2} (l\dot{\theta})^2 = \frac{ml^2\dot{\theta}^2}{2}
    \]
    
    \item \textbf{Потенциальная энергия}:
    \[
    U = mgl(1 - \cos\theta)
    \]
    При малых колебаниях используем приближение \( \cos\theta \approx 1 - \frac{\theta^2}{2} \), откуда
    \[
    U \approx mgl\left(1 - \left(1 - \frac{\theta^2}{2}\right)\right) = \frac{mgl\theta^2}{2}
    \]
\end{itemize}

Таким образом, \textbf{полная энергия} математического маятника:
\[
E = T + U = \frac{ml^2\dot{\theta}^2}{2} + mgl(1 - \cos\theta)
\]
а при малых углах
\[
\boxed{E \approx \frac{ml^2\dot{\theta}^2}{2} + \frac{mgl\theta^2}{2}}
\]

\subsection{Энергия физического маятника}
\textbf{Физический маятник} --- это твёрдое тело произвольной формы, способное совершать колебания вокруг неподвижной горизонтальной оси, не проходящей через его центр масс.

\begin{itemize}
    \item \textbf{Кинетическая энергия:} \( T = \frac{I\omega^2}{2} \), где \( I \) — момент инерции маятника относительно оси вращения, \( \omega = \dot{\theta} \) — угловая скорость.
    \[
    T = \frac{I\dot{\theta}^2}{2}
    \]
    
    \item \textbf{Потенциальная энергия:} Если расстояние от оси вращения до центра масс равно \( L \), то
    \[
    U = mgL(1 - \cos\theta)
    \]
    При малых углах аналогично получаем \( U \approx \frac{mgL\theta^2}{2} \).
\end{itemize}

\textbf{Полная энергия} физического маятника:
\[
E = T + U = \frac{I\dot{\theta}^2}{2} + mgL(1 - \cos\theta)
\]
а при малых колебаниях
\[
\boxed{E \approx \frac{I\dot{\theta}^2}{2} + \frac{mgL\theta^2}{2}}
\]

\subsection{Характеристики волны}
\textbf{Волна} --- это процесс распространения колебаний в пространстве с течением времени. Основные параметры волны:

\begin{itemize}
    \item \textbf{Длина волны ($\lambda$):} расстояние между двумя ближайшими точками, колеблющимися в одинаковой фазе. 
    \item \textbf{Период ($T$):} время одного полного колебания точки среды. 
    \item \textbf{Частота ($\nu$ или $f$):} число полных колебаний в единицу времени:
    \[
    f = 1/T
    \] 
    \item \textbf{Скорость волны ($v$):} скорость распространения фазы колебания:
    \[
    \lambda = v T = \frac{v}{f}
    \]
    \item \textbf{Амплитуда ($A$):} максимальное отклонение точек среды от положения равновесия.
\end{itemize}

Виды волн \textit{по направлению колебаний}:
\begin{itemize}
    \item \textbf{Поперечные волны:} частицы среды колеблются перпендикулярно направлению распространения волны (например, волны в струне, электромагнитные волны).
    \item \textbf{Продольные волны:} частицы среды колеблются вдоль направления распространения волны (например, звук в газах и жидкостях). Представляют собой чередующиеся области сжатия и разрежения.
\end{itemize}

Виды волн \textit{по форме волнового фронта}:
\begin{itemize}
    \item \textbf{Плоские волны:} волновые фронты представляют собой параллельные плоскости.
    \item \textbf{Сферические волны:} источником является точка, волновые фронты — концентрические сферы.
\end{itemize}

\subsection{Эффект Доплера}

\textbf{Эффект Доплера} заключается в изменении частоты волн, регистрируемых приемником, вследствие относительного движения источника и приемника.

Пусть:
\begin{itemize}
    \item $v$ --- скорость звука в среде;
    \item $v_s$ --- скорость источника (source);
    \item $v_r$ --- скорость приемника (receiver);
    \item $f_0$ --- частота, испускаемая источником;
    \item $f$ --- частота, фиксируемая приемником.
\end{itemize}

\subsubsection{Движение приемника при неподвижном источнике}

Если приемник движется навстречу источнику со скоростью $v_r$, то относительная скорость волны относительно приемника увеличивается:
\[
v_{rel} = v + v_r
\]
Длина волны в среде остается неизменной: $\lambda = \frac{v}{f_0}$. Частота, которую зафиксирует приемник:
\[
f = \frac{v_{rel}}{\lambda} = \frac{v + v_r}{v / f_0} = f_0 \left( \frac{v + v_r}{v} \right)
\]

\subsubsection{Движение источника при неподвижном приемнике}

Когда источник движется к приемнику со скоростью $v_s$, он «догоняет» испускаемые им гребни волн. За один период $T = 1/f_0$ источник проходит путь $v_s T$, поэтому эффективная длина волны $\lambda'$ сокращается:
\[
\lambda' = \lambda - v_s T = \frac{v}{f_0} - \frac{v_s}{f_0} = \frac{v - v_s}{f_0}
\]
Приемник фиксирует частоту:
\[
f = \frac{v}{\lambda'} = \frac{v}{(v - v_s)/f_0} = f_0 \left( \frac{v}{v - v_s} \right)
\]

\subsubsection{Общая формула}
Объединяя оба случая (движение вдоль одной прямой), получаем общую формулу эффекта Доплера:
\[
f = f_0 \left( \frac{v \pm v_r}{v \mp v_s} \right)
\]
Где:
\begin{itemize}
    \item Верхние знаки соответствуют \textbf{сближению}.
    \item Нижние знаки соответствуют \textbf{удалению}.
\end{itemize}

\paragraph*{Пример.}
Поезд (\(v_s = 30\) м/с) приближается к наблюдателю. Частота гудка \(f_0 = 500\) Гц, скорость звука \(v = 340\) м/с. Частота, воспринимаемая наблюдателем:
\[
f = f_0 \frac{v}{v-v_s} = 500 \cdot \frac{340}{310} \approx 548 \text{ Гц}
\]
При удалении: \(f = f_0 \frac{v}{v+v_s} = 500 \cdot \frac{340}{370} \approx 460 \text{ Гц}\)

\subsection{Биения}

\textbf{Биения} — это периодические изменения амплитуды результирующего колебания, возникающие при наложении двух гармонических колебаний с близкими частотами.

Рассмотрим сложение двух колебаний одинаковой амплитуды $A$ с близкими частотами $\omega_1$ и $\omega_2$ ($\omega_1 \approx \omega_2$):
\[
x_1(t) = A \cos(\omega_1 t), \quad x_2(t) = A \cos(\omega_2 t)
\]

Согласно принципу суперпозиции, результирующее смещение $x = x_1 + x_2$. Используя тригонометрическую формулу суммы косинусов:
\[
x(t) = 2A \cos\left( \frac{\omega_1 - \omega_2}{2}t \right) \cos\left( \frac{\omega_1 + \omega_2}{2}t \right)
\]

Полученное уравнение можно рассматривать как гармоническое колебание с «быстрой» циклической частотой $\omega_{avg} = \frac{\omega_1 + \omega_2}{2}$ и медленно меняющейся амплитудой.

\begin{itemize}
    \item \textbf{Амплитуда биений} определяется модулем первого множителя:
    \[
    A_{beat}(t) = \left| 2A \cos\left( \frac{\omega_1 - \omega_2}{2}t \right) \right|
    \]
    
    \item \textbf{Частота биений} определяется максимумом интенсивности дважды за период функции косинуса \textit{(в положительном и отрицательном максимуме)} и равна разности исходных циклических частот:
    \[
    \Omega = |\omega_1 - \omega_2|
    \]
    
    \item \textbf{Период биений:} 
    \[
    T_{beat} = \frac{2\pi}{\Omega} = \frac{1}{|f_1 - f_2|}
    \]
\end{itemize}

Биения отчетливо слышны, когда разность частот невелика \textit{(до 10-20 Гц)}. В акустике это воспринимается как периодическое усиление и ослабление звука \textit{(«пульсация»)}.

Настройка инструментов по камертону — классический пример использования биений:
\begin{itemize}
    \item Настройщик одновременно извлекает звук из эталона (камертона) с частотой $f_{std}$ и из настраиваемой струны с частотой $f_{str}$.
    \item Если частоты не совпадают ($f_{std} \neq f_{str}$), слышны характерные пульсации громкости (биения). Громкость звука периодически нарастает и падает с частотой $\Omega = 2\pi|f_{std} - f_{str}|$.
    \item Изменяя натяжение струны, настройщик уменьшает разность частот. Когда пульсации становятся всё более редкими и, наконец, исчезают ($\Omega \to 0$), это означает, что частота струны совпала с эталоном ($f_{str} = f_{std}$), и инструмент настроен «в унисон».
\end{itemize}